\begingroup
\raggedright
\subsection{String Input and getline}
\texttt{cin >> s} reads a whitespace-separated word.
\texttt{getline(cin, s)} reads a whole line, preserving leading and
trailing spaces and accepting empty lines. It consumes the newline
without storing it. Include \texttt{<string>}, \texttt{<iostream>},
\texttt{<sstream>} for parsing, and \texttt{<limits>} for \texttt{ignore}.

\medskip
\textbf{A number followed by lines.} Formatted extraction leaves the
newline in the stream. If the strings start on the next line, discard
the rest of the number's line before calling \texttt{getline}:
\begin{lstlisting}[style=C++,showstringspaces=false]
int n;
string line;
if (cin >> n) {
    cin.ignore(numeric_limits<streamsize>::max(), '\n');
    for (int i = 0; i < n && getline(cin, line); ++i) {
        // Process line; it may contain spaces or be empty.
    }
}
\end{lstlisting}
This intentionally discards any text after the number on that same
line. To read that remainder instead, call \texttt{getline} directly.

\medskip
\textbf{Whitespace and EOF.} \texttt{getline(cin >> ws, s)} is a
shortcut only when all leading whitespace and blank lines may be
skipped. Use plain \texttt{getline} to preserve them.
\texttt{while (getline(cin, s))} reads until EOF, including a final
nonempty line without a trailing newline. The same patterns work
with \texttt{ifstream} for files and \texttt{istringstream} for strings.
\texttt{cin.get(ch)} reads one character including whitespace;
\texttt{cin >> ch} skips whitespace.

\medskip
\textbf{Runnable examples.} The streams below contain sample input;
replace them with \texttt{cin} for judge input. Quoted output makes
spaces and empty strings visible.
\lstinputlisting[style=C++,showstringspaces=false]{../compile/StringInputExamples.cpp}
\endgroup
